\documentclass[11pt]{article}
\usepackage[utf8]{inputenc}
\usepackage{amsmath,amssymb}
\usepackage[margin=1in]{geometry}
\usepackage{hyperref}
\emergencystretch=3em

\title{The multiplicity theorem with arbitrary $(k,h)$, all exponents equal:\\
$Z_d(n)\sim\Big(\prod_ik_i^{-1}\Big)\dfrac{A_{p-1}}{2^{d-1}(d-1)!}\,n^{-p/2}(\log n)^{d-1}$}
\author{Claude, with Daniel Murfet}
\date{2026-09-06}

\begin{document}
\maketitle
\sloppy

\begin{abstract}
The $d$-fold chart standard integral with exponent vectors $k,h:\mathbb N\to\mathbb N$ is defined
recursively, $Z_0(c)=f(c)$, $Z_{d+1}(c)=\int_0^bu^{h_d}Z_d(cu^{k_d})\,du$. When all candidate
exponents $p_i=(h_i+1)/k_i$ coincide at $p$, it reduces exactly to the iterated weighted primitives
of unit~49, $Z_{d+1}(c)=(\prod_{i\le d}k_i^{-1})\,c^{-p}H_d(cb^{\sum k_i})$ with
$H=\texttt{iterDivPrim}(F_{p-1})$, and hence
$Z_{d+2}(\sqrt n)\sim(\prod k_i^{-1})\,\tfrac{A_{p-1}}{2^{d+1}(d+1)!}\,n^{-p/2}(\log n)^{d+1}$:
the multiplicity-$(d+2)$ statement of \texttt{thm:TaylorTree} for arbitrary $(k,h)$. Unit~39 is the
case $k=h+1=1$. The exact reduction was verified by three-dimensional quadrature at $k=(1,2,1)$,
$h=(0,1,0)$ to $2\cdot10^{-16}$. Zero \texttt{sorry}/\texttt{axiom}.
\end{abstract}

\section{Setting}
Each coordinate contributes one chart power substitution (unit~44) turning $u^{h}\,(\cdot)(cu^{k})$
into $\tfrac1k\int_0^{b^k}s^{p-1}(\cdot)(cs)\,ds$, and one weighted scaling. At the innermost level
the kernel becomes $F_{p-1}$; at every further level the factor $(cs)^{-p}$ from the previous level
combines with $s^{p-1}$ to $s^{-1}$, which is exactly the ``divide and integrate'' step of
\texttt{iterDivPrim}, with the scale absorbed into $b^{\sum k_i}$.

\section{The things formalised}
{\small
\begin{verbatim}
noncomputable def iterChartGen (β a b : ℝ) (k h : ℕ → ℕ) : ℕ → ℝ → ℝ
  | 0 => fun c => quadKernel β a c
  | d + 1 => fun c => ∫ u in Ioc 0 b, u ^ (h d) * iterChartGen β a b k h d (c * u ^ (k d))
theorem iterChartGen_succ_eq ...
    (hp : ∀ i, i ≤ d → ((h i : ℝ) + 1) / k i = p) (hc : 0 < c) :
    iterChartGen β a b k h (d + 1) c
      = (∏ i ∈ Finset.range (d + 1), (1 : ℝ) / k i) * c ^ (-p)
        * iterDivPrim (weightedPrimitive β a (p - 1)) d
            (c * b ^ (∑ i ∈ Finset.range (d + 1), k i))
theorem iterDivPrim_isEquivalent (hG : LogBase G A M C δ) (hA : 0 < A)
    (m : ℕ) (hm : 0 < m) :
    (fun L => iterDivPrim G m L) ~[atTop] fun L => A / (m.factorial : ℝ) * Real.log L ^ m
theorem iterChartGen_isEquivalent ... (hp : ∀ i, i ≤ d + 1 → ((h i : ℝ) + 1) / k i = p) :
    (fun n => iterChartGen β a b k h (d + 2) (Real.sqrt n)) ~[atTop]
      fun n => (∏ i ∈ Finset.range (d + 2), (1 : ℝ) / k i)
        * (weightedMass β a (p - 1) / (2 ^ (d + 1) * ((d + 1).factorial : ℝ)))
        * (n ^ (-(p / 2)) * Real.log n ^ (d + 1))
\end{verbatim}
}

\section{Proof strategy}
The reduction is an induction on $d$ generalising the scale $c$. The base case is unit~44's inner
computation; the step rewrites the inner integral by the induction hypothesis at scale $cu^{k_d}$,
factors $(cu^{k})^{-p}=c^{-p}(u^k)^{-p}$, applies the chart power substitution to
$u^{h}\,g(u^{k})$ with $g(s)=s^{-p}H_d(cb^Ss)$, simplifies $s^{p-1}s^{-p}=s^{-1}$, and applies the
scaling lemma with $\gamma=-1$, whose prefactor $(cb^S)^{-(-1+1)}$ is $1$; the result is
\texttt{iterDivPrim\_succ} at the point $cb^{S+k_d}$. The asymptotic of the iterated primitives of
any admissible base follows from unit~49's expansion: the lower-order terms $\sum_{i<m}c_{m,i}\log^iL$
are $o(\log^mL)$ (\texttt{isLittleO\_pow\_pow\_atTop\_of\_lt} composed with $\log$, summed with
\texttt{IsLittleO.sum}) and the remainder is $O(1)=o(\log^mL)$. The final assembly is unit~39's
template with $(\sqrt n)^{-p}=n^{-p/2}$.

\section{Roadblocks and resolutions}
\begin{itemize}
\item The scaling lemma at $\gamma=-1$ leaves the exponent $-(-1+1)$, not $-1+1$; the
\texttt{show} rewrite must match the negated form before \texttt{Real.rpow\_zero}.
\item \texttt{isLittleO\_pow\_pow\_atTop\_of\_lt} lives in \texttt{SpecificAsymptotics}, which was not
in the import closure; import it explicitly.
\item \texttt{IsLittleO.sum} produces a \texttt{Finset}-sum of functions; convert with
\texttt{Finset.sum\_apply} in a \texttt{congr\_left}.
\end{itemize}

\section{What was Mathlib, what was new}
Mathlib: \texttt{isLittleO\_pow\_pow\_atTop\_of\_lt}, \texttt{IsLittleO.sum}, \texttt{tendsto\_pow\_atTop},
\texttt{Finset.prod\_range\_succ}, \texttt{Finset.sum\_range\_succ}, \texttt{Real.mul\_rpow}. New: the
general-$(k,h)$ chart recursion, its exact reduction, and the multiplicity theorem in the all-equal case
with arbitrary exponents.

\section{Lessons}
The abstraction of unit~49 paid off immediately: nothing about the base $F_{p-1}$ had to be reproved,
and the general-$(k,h)$ recursion is the same three-line pointwise identity at every level. Verifying
the exact reduction numerically in three dimensions before writing the induction gave confidence in
the bookkeeping of $\prod k_i^{-1}$ and $b^{\sum k_i}$.

\section{Outcome}
Units~44--50 lift the whole leading-term theory of \texttt{thm:TaylorTree} from $k=h+1=1$ to arbitrary
$(k,h)$: two dimensions in all three regimes, and $d$ dimensions when all exponents coincide. What
remains at this level of the paper is the mixed case in general dimension (some coordinates at the
minimal exponent, others above), whose two-dimensional and $d=3$ instances (units~37, 40, 46) suggest
the answer: the minimal coordinates alone determine the logarithmic power, the others contribute
bounded factors.
\end{document}
